% Work-MAT -- 2. rocnik, hodina 05. Navazujici vypocty inverzi: odmocniny a lomene vyrazy.
\documentclass[11pt,a4paper]{article}
\input{../../sablony/pripravy/priprava.sty}
\nastavPripravu{2. ročník}{Inverzní funkce}{05}{Inverzní funkce III: další předpisy}
\begin{document}
\poznamka{Navazuj na kuchařku z hodiny 04, její kroky neměň. Dnes je cílem zvládnout výpočet u odmocninových a lomených funkcí a vždy správně určit obory. Příklady se mají řešit řádek po řádku, bez vynechávání prohození proměnných. Pro slabší třídu nech racionální lomenou funkci v rozšíření.}

\blok{Kuchařka platí i pro složitější předpisy}
Stejný postup jako u lineárních funkcí používáme i tehdy, pokud je proměnná pod odmocninou nebo ve jmenovateli. Při každé úpravě hlídáme podmínky: odmocnina je nezáporná, jmenovatel nesmí být nulový a obor původní funkce určuje, jaké hodnoty může mít inverze.

\begin{prikladReseniBox}[title={Příklad 1}]
Urči inverzní předpis k $g(x)=\sqrt{x+1}$ s přirozeným definičním oborem. Uveď obory.
\tcblower
\[
\begin{aligned}
g(x)&=\sqrt{x+1} &&\bigm|\ x\to y\\
x&=\sqrt{y+1}\\
x^2&=y+1\\
y&=x^2-1\\
g^{-1}(x)&=x^2-1.
\end{aligned}
\]
$D(g)=\langle-1;+\infty)$, $H(g)=\langle0;+\infty)$, a tedy $D(g^{-1})=\langle0;+\infty)$ a $H(g^{-1})=\langle-1;+\infty)$. Při umocňování musí být $x\ge0$. Kontrola na oboru inverze: $g(g^{-1}(x))=\sqrt{x^2}=|x|=x$ pro $x\ge0$.
\end{prikladReseniBox}
\poznamka{Při srovnání grafů nakresli náčrt podle $y=x$. Žáci často správně spočítají $x^2-1$, ale chybně napíšou definiční obor celé $\mathbb R$.}

\begin{prikladReseniBox}[title={Příklad 2}]
Urči inverzní předpis k $h(x)=\frac{5}{x-2}+1$ a jeho definiční obor.
\tcblower
\[
\begin{aligned}
h(x)&=\frac{5}{x-2}+1 &&\bigm|\ x\to y\\
x&=\frac{5}{y-2}+1\\
x-1&=\frac{5}{y-2}\\
(x-1)(y-2)&=5\\
y-2&=\frac{5}{x-1}\\
y&=2+\frac{5}{x-1}\\
h^{-1}(x)&=2+\frac{5}{x-1}.
\end{aligned}
\]
$D(h)=\mathbb R\setminus\{2\}$, $H(h)=\mathbb R\setminus\{1\}$ a $D(h^{-1})=\mathbb R\setminus\{1\}$. Ověření: $h(3)=6$ a $h^{-1}(6)=3$.
\end{prikladReseniBox}

\blok{Samostatná práce}
\begin{zadaniReseniBox}[title={Samostatná práce 1}]
Urči inverzní funkce podle kuchařky včetně definičních oborů.
\begin{enumerate}[label=\alph*)]
\item $f(x)=\sqrt{x-3}$, $x\ge3$.
\item $g(x)=\frac{3}{x+2}-1$, $x\ne-2$.
\end{enumerate}
\tcblower
\begin{enumerate}[label=\alph*)]
\item
\(\displaystyle\begin{aligned}f(x)&=\sqrt{x-3} &&\bigm|\ x\to y\\x&=\sqrt{y-3}\\x^2&=y-3\\y&=x^2+3\\f^{-1}(x)&=x^2+3\end{aligned}\)

$D(f^{-1})=\langle0;+\infty)$, $H(f^{-1})=\langle3;+\infty)$.
\item
\(\displaystyle\begin{aligned}g(x)&=\frac{3}{x+2}-1 &&\bigm|\ x\to y\\x&=\frac{3}{y+2}-1\\x+1&=\frac{3}{y+2}\\(x+1)(y+2)&=3\\y+2&=\frac{3}{x+1}\\y&=\frac{3}{x+1}-2\\g^{-1}(x)&=\frac{3}{x+1}-2\end{aligned}\)

$D(g^{-1})=\mathbb R\setminus\{-1\}$, $H(g^{-1})=\mathbb R\setminus\{-2\}$.
\end{enumerate}
\end{zadaniReseniBox}

\begin{zadaniReseniBox}[title={Rozšíření 1}]
Urči inverzi k $p(x)=\frac{2x-1}{x+3}$ na $\mathbb R\setminus\{-3\}$ a její definiční obor.
\tcblower
\[
\begin{aligned}
p(x)&=\frac{2x-1}{x+3} &&\bigm|\ x\to y\\
x&=\frac{2y-1}{y+3}\\
x(y+3)&=2y-1\\
xy+3x&=2y-1\\
y(x-2)&=-1-3x\\
y&=\frac{1+3x}{2-x}\\
p^{-1}(x)&=\frac{1+3x}{2-x}.
\end{aligned}
\]
$D(p^{-1})=\mathbb R\setminus\{2\}$. Původní funkce je prostá: po rovnosti $p(a)=p(b)$ a vynásobení přípustnými jmenovateli vyjde $7a=7b$, takže $a=b$.
\end{zadaniReseniBox}

\clearpage
\begin{zadaniReseniBox}[title={Domácí úkol}]
U každé funkce najdi inverzní předpis podle kuchařky. \textbf{Nakresli graf původní i inverzní funkce} do jedné soustavy souřadnic se zrcadlicí přímkou $y=x$. U obou funkcí urči \textbf{definiční obor a obor hodnot}; u lomených funkcí vyznač také asymptoty.
\begin{enumerate}[label=\alph*)]
\item $f(x)=\sqrt{x+5}$.
\item $g(x)=\frac{4}{x-1}+2$.
\item $p(x)=\frac{2x-1}{x+3}$ (lineárně lomená funkce).
\end{enumerate}
\tcblower
\begin{enumerate}[label=\alph*)]
\item
\(\displaystyle\begin{aligned}f(x)&=\sqrt{x+5} &&\bigm|\ x\to y\\x&=\sqrt{y+5}\\x^2&=y+5\\y&=x^2-5\\f^{-1}(x)&=x^2-5\end{aligned}\)

$D(f)=\langle-5;+\infty)$, $H(f)=\langle0;+\infty)$; $D(f^{-1})=\langle0;+\infty)$, $H(f^{-1})=\langle-5;+\infty)$. Pro náčrt: $[-5;0]$, $[-4;1]$, $[-1;2]$ a jejich obrazy $[0;-5]$, $[1;-4]$, $[2;-1]$.
\item
\(\displaystyle\begin{aligned}g(x)&=\frac{4}{x-1}+2 &&\bigm|\ x\to y\\x&=\frac{4}{y-1}+2\\x-2&=\frac{4}{y-1}\\(x-2)(y-1)&=4\\y-1&=\frac{4}{x-2}\\y&=1+\frac{4}{x-2}\\g^{-1}(x)&=1+\frac{4}{x-2}\end{aligned}\)

$D(g)=\mathbb R\setminus\{1\}$, $H(g)=\mathbb R\setminus\{2\}$; $D(g^{-1})=\mathbb R\setminus\{2\}$, $H(g^{-1})=\mathbb R\setminus\{1\}$. Asymptoty $g$: $x=1$, $y=2$; asymptoty $g^{-1}$: $x=2$, $y=1$. Odpovídající body $[0;-2]\leftrightarrow[-2;0]$.
\item
\(\displaystyle\begin{aligned}p(x)&=\frac{2x-1}{x+3} &&\bigm|\ x\to y\\x&=\frac{2y-1}{y+3}\\x(y+3)&=2y-1\\xy+3x&=2y-1\\y(x-2)&=-1-3x\\y&=\frac{1+3x}{2-x}\\p^{-1}(x)&=\frac{1+3x}{2-x}\end{aligned}\)

$D(p)=\mathbb R\setminus\{-3\}$, $H(p)=\mathbb R\setminus\{2\}$; $D(p^{-1})=\mathbb R\setminus\{2\}$, $H(p^{-1})=\mathbb R\setminus\{-3\}$. Asymptoty $p$: $x=-3$, $y=2$; asymptoty $p^{-1}$: $x=2$, $y=-3$. Odpovídající body $[0;-\frac{1}{3}]\leftrightarrow[-\frac{1}{3};0]$.
\end{enumerate}
\end{zadaniReseniBox}
\end{document}
